\documentclass[11pt]{amsart}
\IfFileExists{lmodern.sty}%
  {\usepackage[T1]{fontenc}\usepackage{lmodern}\usepackage{microtype}}%
  {\usepackage[expansion=false]{microtype}}
\usepackage{amssymb,mathtools}
\usepackage{booktabs,array}
\usepackage{enumitem}
\emergencystretch=2.5em
\DeclareMathSizes{10}{10}{7}{6}
\DeclareMathSizes{9}{9}{6}{6}
\DeclareMathSizes{8}{8}{6}{6}
\usepackage[colorlinks=true,linkcolor=blue,citecolor=blue,urlcolor=blue]{hyperref}
\hypersetup{pdftitle={An irrationality measure for zeta(5)},pdfauthor={Sean Litsey}}

% ---------------------------------------------------------------------------
% Location of the Lean formalization.
\newcommand{\FormalizationURL}{https://github.com/slitsey/catalan-irrationality-measure}
\newcommand{\FormalizationLink}{\url{\FormalizationURL}}
% ---------------------------------------------------------------------------

\newtheorem{theorem}{Theorem}[section]
\newtheorem{lemma}[theorem]{Lemma}
\newtheorem{proposition}[theorem]{Proposition}
\newtheorem{corollary}[theorem]{Corollary}
\theoremstyle{definition}
\newtheorem{definition}[theorem]{Definition}
\theoremstyle{remark}
\newtheorem{remark}[theorem]{Remark}
\numberwithin{equation}{section}
\setcounter{tocdepth}{1}

\newcommand{\Z}{\mathbb Z}
\newcommand{\Q}{\mathbb Q}
\newcommand{\R}{\mathbb R}
\newcommand{\C}{\mathbb C}
\newcommand{\F}{\mathbb F}
\newcommand{\N}{\mathbb N}
\newcommand{\lean}[1]{\nolinkurl{#1}}
\newcommand{\arcosh}{\operatorname{arcosh}}
\newcommand{\vG}{v_p^{\mathrm G}}
\newcommand{\cI}{\mathcal I}
\newcommand{\fr}[1]{\{#1\}}
\newcommand{\gref}{g_{\mathrm{par}}}

\begin{document}

\title[An irrationality measure for $\zeta(5)$]{An irrationality measure for $\zeta(5)$}

\author{Sean Litsey}
\email{sean.litsey.math@gmail.com}

\date{November 16, 2026}

\subjclass[2020]{Primary 11J82; Secondary 11J72, 11M06, 15A15, 31A15, 68V20}

\keywords{Irrationality exponent, $\zeta(5)$, Hankel determinants, $p$-adic valuations, logarithmic energy,
equilibrium measures, formal verification, Lean}

\begin{abstract}
Let $\mu$ denote the irrationality exponent. We prove that $|\zeta(5)-p/q|>q^{-87}$ for all integers $p$ and all
sufficiently large integers $q$, so $\mu(\zeta(5))\le87$. Fauzan proved that $\zeta(5)$ is
irrational with rationally normalized Hankel determinants, and his argument gives $\mu(\zeta(5))\le260$.
Giraudin lowered this to $171$ by a sharper bound for a relative matrix norm. We keep Fauzan's determinants and
improve the two other constants of the method. On the arithmetic side a reflection symmetry modulo $p$ of the
reduced Gram form gives extra divisibility of the determinants by the primes $p$ with $K/200\le p\le K/3$, a
second local bound with positive row weights is used for the primes $K/3<p\le K$, and the prime sum is evaluated
exactly over a longer range with a cutoff that grows with $K$. The normalizing factor then grows at most like
$e^{1.288003K^2}$. On the analytic side we bound the logarithmic energy by a comparison measure on one interval,
a Chebyshev series of degree $60$ against the arcsine density whose positivity is an exact sum of squares. Its
potential and energy have closed forms, and the pointwise inequality is checked on $110$ cells in exact
arithmetic. This gives decay at the rate $e^{-1.38500086K^2}$. With Giraudin's form of the relative norm estimate
the exponent is below $86.54$. The proof is formalized in Lean~4.
\end{abstract}

\maketitle
\tableofcontents

%% =====================================================================
\section{Introduction}\label{sec:intro}

\subsection{The result}
The \emph{irrationality exponent} $\mu(\xi)$ of a real number $\xi$ is the infimum of the real numbers $\mu$ such that
$|\xi-p/q|>q^{-\mu}$ for all integers $p$ and all sufficiently large integers $q$. It does not change when $\xi$ is
multiplied by a nonzero rational number or translated by a rational number. A finite bound for $\mu(\xi)$ is an
irrationality measure for $\xi$, and it shows that $\xi$ is irrational. We write
\[
 \zeta(5)=\sum_{n\ge1}\frac1{n^5}=1.03692\,77551\ldots
\]

\begin{theorem}\label{thm:main}
There is an integer $q_0$ such that
\[
 \Bigl|\zeta(5)-\frac pq\Bigr|>q^{-87}
\]
for all integers $p$ and all integers $q\ge q_0$. In particular $\mu(\zeta(5))\le87$.
\end{theorem}

The proof gives the exponent $\lambda\rho/(D-\Lambda)<86.535$, with the constants of Table~\ref{tab:constants}.
Theorem~\ref{thm:main} has been verified in the Lean~4 proof assistant, together with every theorem, proposition and
lemma used in its proof. The formal proof has no unproved steps and uses only the standard axioms.
Section~\ref{sec:lean} describes it. The remarks are not part of the formal proof. The formalization is available at
\FormalizationLink. The use of AI in this work is disclosed at the end of the paper.

\subsection{Background}\label{ssec:background}
\emph{Odd zeta values.} Ap\'ery proved that $\zeta(3)$ is irrational \cite{Apery79}, and Beukers gave his integral
proof \cite{Beukers79}. For the higher odd values the classical results concern families. Rivoal
\cite{Rivoal00} and Ball and Rivoal \cite{BallRivoal01} proved that infinitely many of the numbers $\zeta(2k+1)$ are
irrational. Their bound for the dimension of the $\Q$-span of $1,\zeta(3),\dots,\zeta(s)$ grows like $\log s$.
Zudilin proved that at least one of $\zeta(5)$, $\zeta(7)$, $\zeta(9)$ and $\zeta(11)$ is irrational
\cite{Zudilin01,Zudilin04}. Fischler, Sprang and Zudilin showed that the number of irrational values among
$\zeta(3),\dots,\zeta(s)$ grows faster than any power of $\log s$ \cite{FSZ19}. Fischler improved the dimension bound
to a multiple of $\sqrt{s/\log s}$ with Siegel's lemma \cite{Fischler26}. Fischler, Lai, Sprang and Zudilin raised the
count of irrational values to a multiple of $\sqrt{s\log\log s/\log s}$ with a Hankel form of the construction
\cite{FLSZ26}.

\emph{The number $\zeta(5)$.} Brown and Zudilin constructed simultaneous rational approximations to $\zeta(3)$ and
$\zeta(5)$ from cellular integrals on the moduli space $\mathcal M_{0,8}$. They obtained an effective sequence of
rationals with $0<|\zeta(5)-p/q|<q^{-0.86}$ \cite{BrownZudilin22}. This exponent does not imply irrationality. Fauzan
proved that $\zeta(5)$ is irrational \cite{Fauzan2026}. His approximations are integer polynomials of growing degree,
rationally normalized Hankel determinants of a functional whose value at $\zeta(5)$ is integration against a positive
weight. The same paper proves $\mu(\zeta(5))\le260$ \cite[Corollary~1.2]{Fauzan2026}. Giraudin replaced two
worst-case estimates in Fauzan's bound for a relative matrix norm by closed-form functions of the node, and obtained
$\mu(\zeta(5))\le171$ \cite{Giraudin26z5}. His theorem is stated under the hypotheses of Fauzan's paper. Fasten
\cite{Fasten26}, Firsching \cite{Firsching26} and Romik \cite{Romik26} verified Fauzan's irrationality theorem in Lean.
Romik's formalization takes the prime number theorem as an axiom. As far as we know, $260$ and $171$ are the only
published irrationality measures for $\zeta(5)$. Theorem~\ref{thm:main} improves them to $87$, and its proof is
unconditional and formally verified.

\emph{Recent results.} Since September 2026 several new results have appeared in this area, some obtained with AI
systems. OpenAI proved that Catalan's constant $G$ is irrational \cite{OpenAI2026} and that $\mu(\pi)=2$
\cite{OpenAIPi2026}, both with Lean formalizations. Liu, Jiang and Zhang proved $\mu(\log r)=2$ for rational $r>0$,
$r\ne1$ \cite{LiuJiangZhang26}. Giraudin improved the bounds for $\mu(\zeta(2))$ and $\mu(\zeta(3))$
\cite{Giraudin26,Giraudin26z3}. Lai combined Fauzan's determinants with simultaneous approximations to show that
$1$, $\zeta(3)$ and $\zeta(4)$ are linearly independent \cite{Lai2026}. Leibtag, Barkan, Kalisch, Weinbaum, Shalyt
and Kaminer used conservative matrix fields to improve the measures for $\zeta(2)$ and $\zeta(3)$, and they also
proved this linear independence \cite{LBKWSK26}.

\emph{The method.} Irrationality proofs with determinants of growing size go back to Zudilin's determinantal approach
\cite{Zudilin}. The measure method we use, a family of integer polynomials whose values near the constant are small
and nonzero, is a one-variable form of the criteria of Nesterenko \cite{Nesterenko85} used by Ball and Rivoal
\cite{BallRivoal01}. The arithmetic of the coefficients is handled prime by prime, as in the work of Hata
\cite{Hata93}, Rhin and Viola \cite{RhinViola96,RhinViola01} and Zudilin \cite{Zudilin04}. The size of the
determinants is controlled by Andr\'eief's identity \cite{Forrester2018Andreief} and a logarithmic energy problem with
an external field \cite{SaffTotik}.

\subsection{The companion papers}
This paper applies one method to a third constant. The method consists of Hankel or L\"owner determinants,
denominators bounded prime by prime, decay proved through an energy problem, and a formal proof of every step. The
first companion paper \cite{Litsey2026} refines OpenAI's determinant for $G$ and proves $\mu(G)\le270$. The second
\cite{Litsey2026b} uses L\"owner determinants of digamma combinations and proves measures on three planes of numbers,
spanned modulo $\Q$ by $L(2,\chi_{-3})$ and $\pi^2$, by $\zeta(3)$ and $\pi^2$, and by $\zeta(3)$ and $\pi^4$. We use the definitions of $\mu$, the measure criterion and the
potential-theoretic language of these papers. The construction here is Fauzan's, with his notation.

\subsection{What a measure requires}
Fauzan's irrationality proof needs three facts about his polynomials $Q_n$ of degree $h=37n$: integral coefficients,
$Q_n(\zeta(5))>0$, and $Q_n(\zeta(5))<e^{-cn^2}$. If $\zeta(5)=a/b$, then $b^hQ_n(a/b)$ is a positive integer below
$1$ for large $n$. A measure needs more. Let $p/q$ be a rational close to $\zeta(5)$. Three estimates are needed.
\begin{enumerate}[label=(\roman*)]
\item \emph{Denominators.} A normalizing factor $m_K$ with $m_KF_K\in\Z[X]$ and $\log m_K\le(\Lambda+o(1))K^2$.
\item \emph{Decay on a neighbourhood.} $|F_K(x)|\le e^{-(D-o(1))K^2}$ for every real $x$ with
 $|x-\zeta(5)|\le e^{-(\rho+o(1))K}$, not only at $x=\zeta(5)$.
\item \emph{Nonvanishing near the point.} $F_K(p/q)\ne0$ for every rational $p/q$ in the same neighbourhood.
\end{enumerate}
Then $\mu(\zeta(5))\le\lambda\rho/(D-\Lambda)$, where $\lambda=h/K$ is the degree per unit of scale. For Fauzan's
determinants the matrix is affine in the variable, and it is positive definite at $\zeta(5)$. A relative norm bound
for the coefficient matrix of the variable therefore gives (ii) and (iii) together, with the neighbourhood exponent
$\rho$ of the norm. The constant $\rho$ was bounded by Fauzan \cite[Appendix~C]{Fauzan2026} and sharpened by Giraudin
\cite{Giraudin26z5}. The exponent is very sensitive to the gap $D-\Lambda$. Fauzan's constants give
$D-\Lambda\approx0.0494$, and ours give $D-\Lambda>0.0969$. This halves the exponent.

\subsection{New ideas}
Our proof keeps Fauzan's determinants unchanged and rests on the following ideas.
\begin{enumerate}[leftmargin=*]
\item \emph{The parity gain.} For a prime $p$ in the inner range, Fauzan bounds the $p$-adic valuation of the
 determinant by row weights attached to the classes of poles modulo $p$. On the rows of integral weight we compute the
 leading term of the Gram form modulo $p$. It is a sum over residue classes of local bilinear forms, and each local form
 is odd under the reflection $z\mapsto-1-z$ of its pole set. The classes whose space of local polynomials has
 more even than odd elements under this reflection give radical vectors of the reduced form, and a dimension count
 bounds the loss from the coupling between classes. Each radical vector raises the valuation of the determinant by
 one. For $K/200\le p\le K/3$ this gives $p\,\gref(K/p)+O(M)$ extra
 units, with an explicit density $0\le\gref\le\frac12$ (Theorem~\ref{thm:parity}).
\item \emph{A second bound in the outer range.} For $K/3<p\le K$ we also apply Fauzan's class construction of the
 inner range, with weights that may be positive. The polynomial part of the functional is kept, so the bound holds for
 the whole matrix. Its cost is the inner profile $R$ continued to $1<K/p<3$, and the minimum with Fauzan's outer bound
 saves $3203/163200$ in $\Lambda$ (Theorem~\ref{thm:outer}).
\item \emph{An exact prime sum.} Fauzan's inner profile $R$ has a closed form $R(x)=xF(x)+Q(x)$ with an explicit
 sawtooth $F$ and $-\frac12\le Q\le\frac{13}8$. We integrate $R(x)x^{-3}$ exactly over $[3,200]$ and bound the tail by
 one integration by parts. We let the cutoff between the small and the inner primes grow like $K^{1/4}$. The error terms
 of the local bounds are uniform in the cutoff, so the terms of order $1/M$ of a fixed cutoff disappear
 (Theorem~\ref{thm:lambda}). Together with items 1 and 2 this gives $\Lambda=1.288003$, against $1.31758$.
\item \emph{A comparison measure for the energy problem.} Fauzan bounds the energy with a sum of sixteen arcsine
 measures. We use one measure $\sigma=f\,\omega_{[0,b]}$, where $\omega_{[0,b]}$ is the arcsine probability measure of
 $[0,b]$ and $f$ is a Chebyshev series of degree $60$. Its positivity is an exact sum of squares. Its potential and
 energy have closed forms, inside and outside the support. The pointwise inequality of the energy method is checked on
 $110$ cells by Taylor polynomials of order four with explicit remainders. A layer-cake estimate controls the hard edge
 of $\sigma$ in the regularization. This gives $D=1.38500086$ against $1.36699$ (Theorem~\ref{thm:decay}).
\end{enumerate}
The relative norm exponent $\rho=9.07427$ follows Giraudin's argument \cite{Giraudin26z5} with another choice of
parameters (Section~\ref{sec:rho}). The measure criterion is Fauzan's argument of \cite[Appendix~C]{Fauzan2026},
stated for a family with a rational normalizing factor (Section~\ref{sec:criterion}).

\begin{table}[ht]
\centering
\renewcommand{\arraystretch}{1.2}
\begin{tabular}{@{}lcccc@{}}
\toprule
 & $\Lambda$ & $D$ & $\rho$ & $\mu(\zeta(5))\le$\\
\midrule
Fauzan \cite{Fauzan2026} & $1.31758$ & $1.36699$ & $13.866$ & $260$\\
Giraudin \cite{Giraudin26z5} & $1.31758$ & $1.36699$ & $9.09$ & $171$\\
This paper & $1.288003$ & $1.38500086$ & $9.07427$ & $87$\\
\bottomrule
\end{tabular}
\caption{The three constants of the measure method for Fauzan's determinants, with $\lambda=37/40$. The bound is
$\lambda\rho/(D-\Lambda)$, rounded up. Fauzan's $\Lambda$ is his $A_{100000}$, and his $D$ is $-U$ in
\cite[(6.4)]{Fauzan2026}. His paper and Giraudin's use the margin $D-\Lambda>79/1600$.}\label{tab:constants}
\end{table}

\subsection{Organization}
Section~\ref{sec:construction} recalls Fauzan's construction. Section~\ref{sec:criterion} states the measure criterion
and reduces nonvanishing and decay near $\zeta(5)$ to decay at $\zeta(5)$ and a relative norm bound.
Section~\ref{sec:rho} proves the norm bound. Sections~\ref{sec:local}, \ref{sec:parity} and~\ref{sec:sum} prove the
bound for the normalizing factor: the local bounds, the parity gain and the prime sum. Section~\ref{sec:decay} proves
the decay. Section~\ref{sec:proof} completes the proof. Section~\ref{sec:lean} describes the formalization. The
disclosure of AI use follows it.

\subsection{Notation}
We write $v_p$ for the $p$-adic valuation. For a polynomial $A$ with rational coefficients, $\vG(A)$ is the minimum
of the valuations of its coefficients, the Gauss valuation, with $\vG(0)=+\infty$. We write $\fr y=y-\lfloor y\rfloor$
for the fractional part, and $y_+=\max(y,0)$. We write $\xi=\zeta(5)$ throughout.

%% =====================================================================
\section{Fauzan's construction}\label{sec:construction}

We recall the construction of \cite[\S2]{Fauzan2026}. Throughout,
\begin{equation}\label{eq:params}
 K=40n,\quad N=3n,\quad h=37n,\quad \alpha=\tfrac3{40},\quad\lambda=\tfrac{37}{40},\quad H=1+2\alpha=\tfrac{23}{20}.
\end{equation}
Put $D_m(t)=\prod_{j=1}^m(t+j^2)$ and $H^{(5)}_j=\sum_{v=1}^jv^{-5}$. With $B_1=-\frac12$, the $\Q$-linear functional
$\mu_X$ with values in $\Q[X]$ is defined by
\begin{gather}
 \mu_X(t^e)=\frac{(-1)^eB_{2e+2}(2e+3)(2e+4)(2e+5)}{24}\qquad(e\ge0),\label{eq:functional}\\
 \mu_X\Bigl(\frac1{t+j^2}\Bigr)=j^4\bigl(X-H^{(5)}_j\bigr)-\frac14+\frac1{2j}\qquad(j\ge1),\notag
\end{gather}
and extended to rational functions with simple poles at negative integer squares by polynomial division and partial
fractions. Then
\begin{equation}\label{eq:matrix}
 G_K(X)=\Bigl[\mu_X\Bigl(\frac{D_N(t)^6t^{i+j}}{D_K(t)}\Bigr)\Bigr]_{0\le i,j<h},\qquad \Delta_K=\det G_K,
\end{equation}
\[
 S_K=\frac{(K!)^{2h}4^{h-1}}{(N!)^{12h}\prod_{i=1}^{h-1}((2i)!)^2},\qquad F_K=S_K\Delta_K\in\Q[X].
\]
The entries of $G_K(X)$ are affine in $X$. We write $M_K=[X]G_K(X)$ for the matrix of coefficients of $X$ and
$G=G_K(\xi)$.

\begin{proposition}[Fauzan]\label{prop:basic}
\begin{enumerate}[label=(\alph*)]
\item For every rational function $R$ in the domain of $\mu_X$,
\[
 \mu_\xi(R)=\int_0^\infty R(y^2)\,w(y)\,dy,\qquad w(y)=\frac{(2\pi)^4y^5}{12}\sum_{\ell\ge1}\ell^4e^{-2\pi\ell y}.
\]
\item $G$ is positive definite. In particular $F_K(\xi)>0$.
\item $\deg\Delta_K=h$.
\end{enumerate}
\end{proposition}

Part (a) comes from Euler's formula for $\zeta(2e+2)$ and Hermite's integral for the Hurwitz zeta function
\cite[(2.10)]{Fauzan2026}. Part (b) follows because $w>0$, and part (c) from the Vandermonde factorization
$M_K=V\operatorname{diag}(\dots)V^T$ \cite[(2.9)]{Fauzan2026}. The integral representation also gives Andr\'eief's
formula
\begin{equation}\label{eq:andreief}
 \Delta_K(\xi)=\frac1{h!}\int_{(0,\infty)^h}\prod_{i<j}(y_i^2-y_j^2)^2\prod_{i=1}^h\frac{D_N(y_i^2)^6}{D_K(y_i^2)}\,w(y_i)\,dy .
\end{equation}

%% =====================================================================
\section{The measure criterion}\label{sec:criterion}

\subsection{The deduction}
The criterion is that of the companion papers \cite{Litsey2026,Litsey2026b}, in the form with nonvanishing near the
point. Let $(P_n)$ be polynomials with rational coefficients, indexed by $n\ge n_0$, at the scale $K=40n$.

\begin{theorem}[measure criterion]\label{thm:criterion}
Let $\xi$ be real and let $\Lambda<D$, $\rho>0$ and $\gamma>0$ be real numbers. Suppose that $\deg P_n\le\gamma K$, and
that for every $\varepsilon>0$ and all large $n$ the following hold.
\begin{enumerate}[label=(\alph*)]
\item There is a positive integer $d_n$ with $\log d_n\le(\Lambda+\varepsilon)K^2$ and $q^{\deg P_n}d_nP_n(p/q)\in\Z$
 for all integers $p$ and $q>0$.
\item $|P_n(\xi+t)|\le e^{-(D-\varepsilon)K^2}$ for every real $t$ with $|t|\le e^{-(\rho+\varepsilon)K}$.
\item $P_n(x)\ne0$ for every rational $x$ with $|\xi-x|\le e^{-(\rho+\varepsilon)K}$.
\end{enumerate}
Then $\xi$ is irrational. Moreover, for every $\mu_0>\gamma\rho/(D-\Lambda)$ there is an integer $q_0$ such that
$|\xi-p/q|>q^{-\mu_0}$ for all coprime integers $p$ and $q$ with $q\ge q_0$.
\end{theorem}

\begin{proof}
Let $\eta=D-\Lambda$. If $\xi=a/b$ were rational, with $b>0$, then $\xi$ lies in every window of (c). So
$b^{\deg P_n}d_nP_n(\xi)$ is a nonzero integer. By (a) and (b) with $t=0$ its absolute value is at most
$b^{\gamma K}e^{-(\eta-2\varepsilon)K^2}$, which is below $1$ for small $\varepsilon$ and large $n$. This is
impossible, so $\xi$ is irrational.

Now let $\mu_0>\mu_1>\gamma\rho/\eta$, and suppose $|\xi-p/q|\le q^{-\mu_0}$ with $q$ large. Take the least $n$ with
$\rho K\ge\mu_1\log q$. Then $\rho K\le\mu_1\log q+40\rho$, so $|\xi-p/q|\le q^{-\mu_0}\le e^{-(\rho+\varepsilon)K}$
for a small fixed $\varepsilon$ and large $q$. By (c), $z=q^{\deg P_n}d_nP_n(p/q)$ is a nonzero integer, so $|z|\ge1$.
By (a) and (b), $|z|\le q^{\gamma K}e^{(\Lambda+\varepsilon)K^2-(D-\varepsilon)K^2}$. Hence
$(\eta-2\varepsilon)K\le\gamma\log q$. But $\eta K\ge(\mu_1\eta/\rho)\log q$ and $\mu_1\eta/(\gamma\rho)>1$. This is a
contradiction for small $\varepsilon$ and large $q$.
\end{proof}

The second statement concerns fractions in lowest terms. The following step, used also in the companion papers,
removes this restriction once $\xi$ is known to be irrational.

\begin{lemma}\label{lem:plain}
Let $\xi$ be irrational and $E\ge0$. Suppose that for every $e>0$ there is $q_0$ with $|\xi-p/q|>q^{-(E+e)}$ for all
coprime $p$ and $q$ with $q\ge q_0$. Then for every $\mu_0>E$ there is $q_1$ with $|\xi-p/q|>q^{-\mu_0}$ for all
integers $p$ and all integers $q\ge q_1$.
\end{lemma}

\begin{proof}
Take $e=\mu_0-E$ and the corresponding $q_0$. Write $p/q=p'/q'$ in lowest terms, so $q'\le q$. If $q'\ge q_0$, then
$|\xi-p/q|>q'^{-\mu_0}\ge q^{-\mu_0}$. The fractions with $q'<q_0$ have finitely many denominators. Since $\xi$ is
irrational, $c=\min_{q'<q_0}\min_{p'}|\xi-p'/q'|$ is positive. For $q>c^{-1/\mu_0}$ we get $|\xi-p/q|\ge c>q^{-\mu_0}$.
\end{proof}

Fauzan's normalizing factor $m_K$ is a positive rational number with exponents of both signs. A positive exponent
divides $F_K$ by a power of $p$, and a negative one clears a denominator. So $m_K$ is not an integer multiple of a
denominator. We apply Theorem~\ref{thm:criterion} to the family $P_n=m_KF_K$ itself. Its coefficients are integers, so
(a) holds with $d_n=1$ and $\Lambda=0$. If $\log m_K\le(\Lambda+o(1))K^2$ and $F_K$ satisfies (b) with the rate $D$,
then $P_n$ satisfies (b) with the rate $D-\Lambda$. The exponent is the same.

\subsection{Nearby nonvanishing from a relative norm bound}
The following lemma is the version of \cite[Appendix~C]{Fauzan2026} that we use.

\begin{lemma}[relative norm]\label{lem:relnorm}
Let $G$ be a real positive definite $h\times h$ matrix and $M$ a real symmetric $h\times h$ matrix with
$|v^TMv|\le\beta\,v^TGv$ for all $v\in\R^h$. If $|\delta|\beta h\le\frac12$, then $0<\det(G+\delta M)\le2\det G$.
\end{lemma}

\begin{proof}
Write $G+\delta M=G^{1/2}(I+\delta B)G^{1/2}$ with $B=G^{-1/2}MG^{-1/2}$. The eigenvalues $\beta_i$ of $B$ satisfy
$|\beta_i|\le\beta$. So $1+\delta\beta_i\ge\frac12>0$, and $\prod_i(1+\delta\beta_i)\le e^{|\delta|\beta h}\le e^{1/2}\le2$.
\end{proof}

\begin{definition}
We say that $\rho$ is a \emph{norm exponent} if $|v^TM_Kv|\le e^{\rho K}\,v^TGv$ for all $v\in\R^h$ and all large $n$.
\end{definition}

Since $G_K(x)=G+(x-\xi)M_K$, Lemma~\ref{lem:relnorm} gives the following.

\begin{corollary}\label{cor:nearby}
Let $\rho$ be a norm exponent. If $|x-\xi|e^{\rho K}h\le\frac12$, then $0<F_K(x)\le2F_K(\xi)$. Hence for every
$\varepsilon>0$ and all large $n$, $F_K(x)\ne0$ and $|F_K(x)|\le2F_K(\xi)$ whenever $|x-\xi|\le e^{-(\rho+\varepsilon)K}$.
\end{corollary}

The second statement uses $2he^{-\varepsilon K}\le1$ for large $K$. So the decay at $\xi$ alone gives (b), and
positivity gives (c).

%% =====================================================================
\section{The norm exponent}\label{sec:rho}

We follow Giraudin \cite{Giraudin26z5}, who refines \cite[Proposition~C.1]{Fauzan2026}. The field of the energy problem
is
\begin{equation}\label{eq:field}
 V(t)=2\pi\sqrt t+\int_0^1\log(t+u^2)\,du-6\int_0^\alpha\log(t+u^2)\,du\qquad(t\ge0).
\end{equation}
Let $g(x)=x(\log|x|-1)$ with $g(0)=0$, and for $\alpha<u\le1$ put
\[
 L(u)=6\bigl[g(\alpha-u)+g(\alpha+u)\bigr]-g(1-u)-g(1+u).
\]
For $0<a<b$ put $Y(u)=(2u^2+a+b)/(b-a)$ and $\Phi_{a,b}(u)=L(u)+2\lambda\arcosh Y(u)$.

\begin{proposition}\label{prop:rho}
The number $\rho=907427/10^5$ is a norm exponent.
\end{proposition}

\begin{proof}
We list the steps. Write $u_j=j/K$, let $q\in\R[t]$ have degree less than $h$, and put $p(t)=q(K^2t)$.
\begin{enumerate}[label=(\arabic*),leftmargin=*]
\item \emph{Residues.} $q^TM_Kq=A_K\sum_{N<j\le K}r_jp(-u_j^2)^2$ with $A_K=K^{12N-2K+6}$ and
 $r_j=u_j^4\prod_{l\le N}(u_l^2-u_j^2)^5/\prod_{N<k\le K,k\ne j}(u_k^2-u_j^2)$. This is the Vandermonde factorization of
 $M_K$.
\item \emph{Nodes.} $|r_j|\le\exp(KL(u_j))$. The logarithm of $|r_j|$ is a Riemann sum of the monotone functions
 $\log|u^2-v^2|$, compared with the antiderivatives $g(u+v)\mp g(u-v)$ by the mean value theorem.
\item \emph{Gram matrix.} For $W\ge\max_{[a,b]}V$,
 $q^TGq\ge A_Kc_{a,b}e^{-KW}\int_a^bp(t)^2\,dt$ with $c_{a,b}=(2\pi)^4a^3/(24(1+b))$. Keep the term $\ell=1$ of $w$,
 restrict to $y=K\sqrt t$ with $a\le t\le b$, and compare the products $D_N^6/D_K$ with $e^{-KV}$ by Riemann sums.
\item \emph{Legendre polynomials.} If $\deg P<h$ and $u>0$, then
 $P(-u^2)^2\le h^2e^{2h\arcosh Y(u)}(b-a)^{-1}\int_a^bP^2$. Expand $P$ in the Legendre polynomials of $[a,b]$, use
 Cauchy--Schwarz, and bound the Legendre polynomials outside the interval through Rodrigues' formula.
\item \emph{Assembly.} Steps (1)--(4) give
 $|q^TM_Kq|\le h^3((b-a)c_{a,b})^{-1}\exp\bigl(K(W+\sup_{\alpha<u\le1}\Phi_{a,b})\bigr)\,q^TGq$.
\item \emph{Shape of $V$.} With $y=\sqrt t$, the derivative of $V(y^2)$ is $2(\pi+\arctan(1/y)-6\arctan(\alpha/y))$.
 This is increasing on $(0,\frac3{20}]$ and positive on $[\frac3{20},\infty)$. So $\max_{[a,b]}V=\max(V(a),V(b))$.
\item \emph{Concavity.} If $ab\le\alpha^4$, then $\Phi_{a,b}'$ is decreasing on $(\alpha,1)$. So
 $\Phi_{a,b}(u)\le\Phi_{a,b}(u_0)+\Phi_{a,b}'(u_0)(u-u_0)$ for every $u_0\in(\alpha,1)$.
\item \emph{Numbers.} Take $a=10^{-8}$, $b=242976/10^6$ and $u_0=973660012/10^9$. Then $V(a)$ and $V(b)$ are at most $W=30449413/10^7$, $\Phi_{a,b}(u_0)\le60293244/10^7$ and
 $|\Phi'_{a,b}(u_0)|\le10^{-7}$. By (7), $\sup_{\alpha<u\le1}\Phi_{a,b}\le60293245/10^7$. Hence
 $W+\sup\Phi_{a,b}\le9.0742658<\rho$, and the margin absorbs the factor $h^3((b-a)c_{a,b})^{-1}$ for large $K$.
\end{enumerate}
The numbers in (8) are rational enclosures of finitely many logarithms and arctangents.
\end{proof}

Steps (1)--(6) are Giraudin's Lemmas~2, 3 and~5 and his Proposition~7 \cite{Giraudin26z5}. He takes
$a=10^{-3}$, $b=49/200$ and obtains $9.09$. He locates the maximum of $\Phi_{a,b}$ by its unique critical point. We use
the tangent bound of step (7) instead, which needs one evaluation of $\Phi_{a,b}$ and of its derivative.

%% =====================================================================
\section{Local bounds for the denominators}\label{sec:local}

\subsection{Fauzan's bounds}
Let $p$ be a prime. A lower bound $\vG(F_K)\ge L_p$ allows the factor $p^{-L_p}$ in the normalization. If $L_p<0$ this
clears a denominator, and if $L_p>0$ it removes a power of $p$ common to all coefficients. Fauzan proves such bounds in
three ranges \cite[\S\S3--5]{Fauzan2026}. Fix a cutoff $M\ge40$ with $K\ge200M^2$.

\emph{Small primes}, $pM\le K$. Here $\vG(F_K)\ge-6h\lfloor\log_p5K\rfloor-hv_p(24)$
\cite[(3.12)]{Fauzan2026}.

\emph{Inner primes}, $K/M<p\le K/3$. Let $m=(p-1)/2$. For $0\le a\le m$ and $A\ge1$ let
$\ell_A(a)=\#\{1\le j\le A:j\equiv\pm a\pmod p\}$, and put $b_a=3\ell_N(a)$. With the reserve $L_0=4M+10$ define
integers $T$ and $E$ by $mT+E=h-L_0+3(N-\lfloor N/p\rfloor)$ and $0\le E<m$. The value $\epsilon_a=1$ is given to $E$
classes $a\ne0$ of largest $\ell_K(a)$, and $\epsilon_a=0$ to the others. Put $L_a=T-b_a+\epsilon_a$. The rows of the
basis are
\begin{equation}\label{eq:basis}
 E_{a,i}(t)=\prod_{c\ne a}(t+c^2)^{L_c}\,(t+a^2)^i\qquad(0\le a\le m,\ 0\le i<L_a),
\end{equation}
and for $a\ne0$ the row $(a,i)$ has the weight
\begin{equation}\label{eq:weights}
 w_{a,i}=i+b_a-\frac{\ell_K(a)+4}2 .
\end{equation}
The zero class has the weights of \cite[(4.7)]{Fauzan2026}. The basis is unimodular over $\Z_{(p)}$, the entry in rows
$s$ and $t$ has Gauss valuation at least $w_s+w_t$, and so $\vG(\Delta_K)\ge\gamma^{\rm in}_p=2\sum_sw_s$
\cite[Proposition~4.1]{Fauzan2026}.

\emph{Outer primes}, $K/3<p\le K$. Fauzan separates the residue terms from a polynomial correction of bounded rank.
With nonpositive half-integer weights and a complementary-minor lemma he obtains an explicit integer $\gamma^{\rm out}_p$
with $\vG(\Delta_K)\ge\gamma^{\rm out}_p$ \cite[(4.14) and Proposition~4.3]{Fauzan2026}. For $p>K$ the matrix is
integral, and for $p>2h$ the factor $S_K$ is a $p$-adic unit.

\subsection{The outer extension}
Fauzan uses the class basis \eqref{eq:basis} only in the inner range. It can be used in the outer range too. There
the weights \eqref{eq:weights} are not clipped at zero, and the polynomial part of
the functional must be kept in the local bounds.

\begin{theorem}[outer extension]\label{thm:outer}
Let $p\ge101$ be a prime with $K/3<p\le K$. Reserve $L_0=8$ rows at the zero class, and define $T$, $E$ by
$mT+E=h-8+3N$ with $0\le E<m$. Give the $E$ extra rows to the classes of largest $2\ell_K(a)+[a\le N]$, and put
$L_a=T-3[a\le N]+\epsilon_a$. With the weights \eqref{eq:weights} on the classes $a\ne0$,
\[
 \vG(\Delta_K)\ge\gamma^{\rm ext}_p=2\sum_sw_s ,
\]
and $\gamma^{\rm ext}_p$ is an integer. Moreover $|\gamma^{\rm ext}_p-p\,\Gamma(K/p)|\le83$, where $\Gamma$ is Fauzan's
inner function \cite[(5.4)]{Fauzan2026}.
\end{theorem}

\begin{proof}[Proof sketch]
Since $p\ge101$ and $p\le K$ we have $n\ge3$. Since $p>2N$, $\ell_N(a)=[a\le N]$ for $a\ne0$, so $b_a=3[a\le N]$.
Counting rows gives $2\le T\le6$ and $L_0+\sum_aL_a=h$. If $T=2$, every class with $a\le N$ receives an extra row,
because the key $2\ell_K(a)+[a\le N]$ puts these classes first among those of equal $\ell_K(a)$. Hence $L_a\ge0$ for
every $a$. We then apply the class bound of the inner range \cite[Proposition~4.1]{Fauzan2026} to the whole local
rational function, polynomial part included. It needs three conditions. The first is separation and degree,
$2K<p^2$ and $104n+6<p^2$, which hold because $p>K/3$. The second is that the top weights of two classes $a,c\ne0$
differ by at most $1$, which holds because $\ell_K$ takes two consecutive values and the extra rows go to the larger
ones. The third compares the top weight of each class $a\ne0$, minus one, with the zero class. It holds because
$T\le6$, $\lfloor K/p\rfloor\le2$ and $L_0=8$. The approximation by $p\Gamma(K/p)$ is the count of classes by their
pole numbers, which are floor functions of $K/p$ and $a/p$.
\end{proof}

\subsection{The normalizing factor}
Fauzan's inner profile is $R(x)=-\Gamma(x)-\mathcal N(x)$ \cite[(5.6)]{Fauzan2026}, where
$\mathcal N(x)=2\lambda x\lfloor x\rfloor-12\lambda x\lfloor\alpha x\rfloor-2J(\lambda x)$ and
$J(u)=\lfloor2u\rfloor u-\lfloor2u\rfloor(\lfloor2u\rfloor+1)/4$. It has a closed form. Put $\tau=\fr{2Hx}$,
$s=\fr{2x}$, $\vartheta=\fr{2\lambda x}$, $t_e=\min(\fr x,1-\fr x)$, $t_n=\min(\fr{\alpha x},1-\fr{\alpha x})$, and let
$\epsilon_f$ and $\epsilon_g$ be $+1$ or $-1$ according as $\fr x<\frac12$ and $\fr{\alpha x}<\frac12$.

\begin{lemma}[closed profile]\label{lem:profile}
$R(x)=xF(x)+Q(x)$, where $F(x)=4\lambda+2\lambda\fr x-12\lambda\fr{\alpha x}$ and
\begin{multline*}
 Q(x)=\tfrac12\bigl[\tau(\tau-s)-(\tau-s)_++\vartheta(1-\vartheta)\bigr]\\
 +9t_n(1-2t_n)-3\epsilon_f\epsilon_g\bigl(\min(t_e,t_n)-2t_et_n\bigr).
\end{multline*}
Moreover $-\frac12\le Q\le\frac{13}8$.
\end{lemma}

The proof is a direct evaluation of the integral in $\Gamma$, cell by cell in the fractional parts. The scalar factor
satisfies $0\le v_p(S_K)-p\mathcal N(K/p)\le\lfloor2h/p\rfloor$ whenever $p^2>2K$ (Legendre's formula). Fauzan's inner
estimate becomes uniform in the cutoff:
\begin{equation}\label{eq:innerunif}
 \bigl|-v_p(S_K)-\gamma^{\rm in}_p-p\,R(K/p)\bigr|\le7501M^2\qquad(K/M<p\le K/3).
\end{equation}
In the outer range $|(-v_p(S_K)-\gamma^{\rm ext}_p)-p\,R(K/p)|\le86$. For $K/3<p\le2h$, Fauzan's outer cost is
$-v_p(S_K)-[p\le K]\,\gamma^{\rm out}_p=K\,T_0(p/K)+O(1)$, with error at most $10$, where
\[
 T_0(y)=R_0(y)-d(y)-2\lambda\lfloor1/y\rfloor+\sum_{j=1}^5(2\lambda-jy)_+
\]
is the integrand of \cite[(5.10)]{Fauzan2026}, with $R_0$ and $d$ as in \cite[(5.8), (5.9)]{Fauzan2026}.

\begin{definition}[the local exponents]\label{def:LpC}
Let $M\ge40$ and $K\ge200M^2$. For a prime $p\le2h$ put
\[
 L_p=\begin{cases}
 -6h\lfloor\log_p5K\rfloor-hv_p(24), & pM\le K,\\
 v_p(S_K)+\gamma^{\rm in}_p+[K\le200p]\,k_p, & K/M<p\le K/3,\\
 v_p(S_K)+\max(\gamma^{\rm out}_p,\gamma^{\rm ext}_p), & K/3<p\le K,\\
 v_p(S_K), & K<p\le2h,
 \end{cases}
\]
where $k_p$ is the parity gain of Section~\ref{sec:parity}. Put $m_K=\prod_{p\le2h}p^{-L_p}$. Since
$K\ge200M^2\ge320000$, every prime in the third case satisfies $p>K/3>101$, so Theorem~\ref{thm:outer} applies.
\end{definition}

\begin{proposition}\label{prop:normaliser}
$m_K$ is a positive rational number and $m_KF_K\in\Z[X]$. Also $\log m_K=-\sum_{p\le2h}L_p\log p$.
\end{proposition}

\begin{proof}
At each prime $p\le2h$ the bound $\vG(F_K)\ge L_p$ holds. For the small primes it is Fauzan's \cite[(3.12)]{Fauzan2026}.
In the inner range it is \cite[Proposition~4.1]{Fauzan2026}, and Theorem~\ref{thm:parity} below when $K\le200p$. In the outer range it is Fauzan's
\cite[Proposition~4.3]{Fauzan2026} or Theorem~\ref{thm:outer}, whichever is larger. For $p>K$ it is
\cite[Proposition~4.3]{Fauzan2026}. At $p>2h$ both $S_K$ and $\Delta_K$ are $p$-integral.
\end{proof}

%% =====================================================================
\section{The parity gain}\label{sec:parity}

\subsection{Statement}
Let $p$ be an inner prime, $K/M<p\le K/3$, with the allocation of Section~\ref{sec:local} and $L_0=4M+10$. The
gain is used only when $K\le200p$. By
\eqref{eq:weights}, $w_{a,i}$ is an integer exactly when $\ell_K(a)$ is even. The counts $\ell_K(a)$ for $a\ne0$ take
two consecutive values. Let $\mathcal S=\{a\ne0:\ell_K(a)\text{ even}\}$ and let $\mathcal E=\{a:\epsilon_a=1\}$. Let $\mathcal S_0$ and $\mathcal S_1$ be the sets of $a\in\mathcal S$ with
$\ell_N(a)$ even and odd, respectively. Define
\[
 k_p=\begin{cases}
 \bigl[\#(\mathcal S_0\setminus\mathcal E)-\#\{b\in\mathcal S_1\cap\mathcal E:L_b>0\}-L_0\bigr]_+ ,& T\text{ odd},\\[2pt]
 \bigl[\#(\mathcal S_0\cap\mathcal E)-\#(\mathcal S_1\setminus\mathcal E)-L_0\bigr]_+ ,& T\text{ even}.
 \end{cases}
\]

\begin{theorem}[parity gain]\label{thm:parity}
For every inner prime $p$ with $K\le200p$, $\vG(\Delta_K)\ge\gamma^{\rm in}_p+k_p$.
\end{theorem}

The count $k_p$ has a density. Put $t_e$, $t_n$ as in Lemma~\ref{lem:profile}, $s'=\frac12-t_e$,
$r_n=(t_n-t_e)_+$, $T(x)=\lfloor\frac{23}{10}x\rfloor$ and $u=\frac{23}{20}x-\frac12T(x)$. Let $e=(u-t_e)_+$ if
$\fr x<\frac12$ and $e=\min(u,s')$ otherwise. Define
\[
 \gref(x)=\begin{cases}(e-r_n)_+, & T(x)\text{ even},\\ (s'-e-r_n)_+, & T(x)\text{ odd}.\end{cases}
\]
Then $0\le\gref\le\frac12$, and $\gref$ is periodic with period $40$.

\begin{proposition}\label{prop:kcount}
For every inner prime $p$, $k_p\ge p\,\gref(K/p)-30M$.
\end{proposition}

The proof counts the classes $a\in[1,m]$ by the values of $\ell_K(a)$, $\ell_N(a)$ and $\epsilon_a$. These depend on
$a/p$ through floor functions of $K/p\pm a/p$ and $N/p\pm a/p$, and each type of class fills an interval of $a/p$ up to
$O(1)$ points. The reserve $L_0$ and the boundary effects cost at most $30M$.

\subsection{The leading term modulo $p$}
We now describe the proof of Theorem~\ref{thm:parity}. Following Fauzan \cite[\S3]{Fauzan2026}, the substitution
$t=-x^2$ turns $\mu_X$ into a functional $\tau_X$ on rational functions of $x$ with simple integer poles,
$\mu_X(R)=\tau_X(x^5R(-x^2))$. On polynomials $\tau(P)=\mathcal L(P''')/24$, where $\mathcal L(x^k)=B_k$, and
$\tau_X(1/(x-r))=H^{(5)}_{d(r)}-X$ with $d(r)=r$ for $r\ge0$ and $d(r)=-r-1$ for $r<0$.

We first replace the factors $t+c^2$ of \eqref{eq:basis} by the centred factors $t+c^2-cp$. Under $x=c+pz$ the centred
factor becomes $p(-2c(z+\frac12)-pz^2)$, with leading part proportional to $z+\frac12$. The centred basis is again
unimodular, and its entries satisfy the same bounds $w_s+w_t$. So Fauzan's bound $\vG(\Delta_K)\ge\gamma^{\rm in}_p$
holds for it, with the same weights.

For a finite set $I\subset\Z$ put $B_I(z)=\prod_{k\in I}(z-k)$. For a residue $c$ modulo $p$ let $I_c$ be the set of
$k$ with $c+pk$ a pole, the local pole set. For a polynomial $W$ over $\F_p$ define the local symbol
\begin{equation}\label{eq:localform}
 C_I(W)=\tau\Bigl(\Bigl\lfloor\frac{W}{B_I}\Bigr\rfloor\Bigr)
 +\sum_{k\in I}\operatorname{res}_{z=k}\frac{W(z)}{B_I(z)}\,H^{(5)}_{d(k)}\pmod p,
\end{equation}
where $\lfloor\cdot\rfloor$ is the polynomial part, and the bilinear form $C_I(U,V)=C_I(UV)$. Both parts are needed. The polynomial part has degree up to about
$p$, and it contributes at the same level as the residues.

\begin{lemma}[leading term of $\tau_X$]\label{lem:taulead}
Let $A\in\Q[x]$ and let $\Pi$ be a product of distinct linear factors $x-r$, with the poles distinct modulo $p$ inside
each residue class and the degrees below $p^2$. Suppose that $A(c+pz)$ has Gauss valuation at least $e(c)$ for every
$c$, and put $\beta=\min_c(e(c)-\#\{r\equiv c\})$. Then $p^{4-\beta}\tau_X(A/\Pi)$ has integral constant term, and its
reduction modulo $p$ is the sum over the residues $c$ with $e(c)-\#\{r\equiv c\}=\beta$ of $u_c^{-1}C_{I_c}(\bar a_c)$.
Here $\bar a_c$ is the reduction of $p^{-e(c)}A(c+pz)$ and $u_c$ is the product of $c-r$ over the poles $r\not\equiv c$.
The coefficient of $X$ in $p^{4-\beta}\tau_X(A/\Pi)$ has valuation at least $5$.
\end{lemma}

The proof has three ingredients. The distribution identity $\tau(P)=p^{-4}\sum_{c<p}\tau(P(c+pz))$ follows from the
difference equation $\tau(P(x+1))-\tau(P)=P''''(0)/24$ and Newton's expansion. The Taylor coefficients of the binomial
polynomials at $c+pz$ of degree above $p$ are divisible by $p^2$, and the Tate-algebra bound of
\cite[Lemma~3.1]{Fauzan2026} disposes of them. Finally $p^5H^{(5)}_{d(c+pk)}\equiv H^{(5)}_{d(k)}\pmod{p^5}$, since the
multiples of $p$ up to $d(c+pk)$ are the $p\,j$ with $j\le d(k)$.

\subsection{Reflection and radical vectors}
Write $J$ for the reflection $W(z)\mapsto W(-1-z)$. If $I$ is symmetric under $k\mapsto-1-k$ and $\#I$ is even, then
$B_I$ is $J$-invariant and
\begin{equation}\label{eq:odd}
 C_I(JU,JV)=-C_I(U,V).
\end{equation}
This comes from the reflection identity $\tau_X(g(-1-x))=-\tau_X(g(x))$ \cite[(3.2)]{Fauzan2026} and from
$H^{(5)}_{d(-1-k)}=H^{(5)}_{d(k)}$. In particular $C_I(U,V)=0$ when $U$ and $V$ are both $J$-invariant.

Normalize the Gram matrix by $\tilde G_{st}=p^{-\lfloor w_s\rfloor-\lfloor w_t\rfloor}G_{st}$. For rows $s$ and $t$ of
integral weight, Lemma~\ref{lem:taulead} identifies $\tilde G_{st}\bmod p$. The residues at the minimal level lie in
the classes of $\mathcal S$, and for such a class the local pole set has $\ell_K(a)$ elements, an even number, and is
symmetric. The residues $c$ and $p-c$ give equal contributions. The result is
\begin{equation}\label{eq:leadform}
 \tilde G_{st}\equiv\kappa_s\kappa_t\sum_{a\in\mathcal S}\delta_a\,C_{I_a}(\Psi_a e_s,\Psi_a e_t)\pmod p .
\end{equation}
Here $\kappa_{a,i}=(-2a)^i$, the $\delta_a$ are units, and $e_s$ is the unit vector of the row $s$. The variable $z$ is
the local variable shifted by $\frac12$, so that $J$ acts by $z\mapsto-z$. The polynomial $B_a$ is the cube of the
local pole polynomial of $D_N$, and it is $J$-invariant when $\ell_N(a)$ is even. Call a class $b\in\mathcal E$ with
$L_b>0$ active, and write $t_b(v)=v_{b,L_b-1}$ for the coordinate of its top row. For a vector $v$ on the rows of the
classes in $\mathcal S$, the local image at $a$ is
\[
 \Psi_av=B_a\Bigl(\sum_{i<L_a}v_{a,i}z^i+[a\notin\mathcal E]\,s_a(v)\,z^{L_a}\Bigr),\qquad
 s_a(v)=\sum_{b\text{ active}}c_{ab}\,t_b(v),
\]
with constants $c_{ab}\in\F_p$. So a row of another class $b$ reaches the class $a$ only through the top row of an
active class $b$, and only when $a$ is not extra.

We now count radical vectors. For $a\in\mathcal S_0$ we have $L_a\equiv T+\epsilon_a\pmod2$. Let $V$ be the space of
vectors supported on the rows $(a,i)$ with $a\in\mathcal S_0$ and $i$ even. Its dimension is the number of rows
$(a,i)$ with $a\in\mathcal S_0$ and $i$ odd, plus the number of $a\in\mathcal S_0$ with $L_a$ odd. By \eqref{eq:odd},
the form \eqref{eq:leadform} vanishes between two vectors whose local images are $J$-invariant at every class of
$\mathcal S_0$ and zero at every class of $\mathcal S_1$.

Let $T$ be odd. Then $L_a$ is odd for $a\in\mathcal S_0\setminus\mathcal E$ and even for $a\in\mathcal S_0\cap\mathcal E$.
The top row of a class in $\mathcal S_0\cap\mathcal E$ has odd index, so $t_b$ vanishes on $V$ for every active $b$
in $\mathcal S_0$, and also for every $b\in\mathcal S_1$. Hence $\Psi_av$ is $J$-invariant for $a\in\mathcal S_0$ and
zero for $a\in\mathcal S_1$, for every $v\in V$. The form between $v\in V$ and $e_r$ can then be nonzero only when $r$
is a row $(a,i)$ with $a\in\mathcal S_0$ and $i$ odd, or the top row of an active class in $\mathcal S_1$. These are
the only linear conditions. So the radical contains at least
$\#(\mathcal S_0\setminus\mathcal E)-\#\{b\in\mathcal S_1\cap\mathcal E:L_b>0\}$ independent vectors of $V$.

Let $T$ be even. Then $L_a$ is odd exactly for $a\in\mathcal S_0\cap\mathcal E$, and for $a\in\mathcal S_0\setminus
\mathcal E$ the power $z^{L_a}$ is $J$-invariant. So $\Psi_av$ is $J$-invariant for $a\in\mathcal S_0$. For
$a\in\mathcal S_1\cap\mathcal E$ we have $\Psi_av=0$. For $a\in\mathcal S_1\setminus\mathcal E$ we impose the condition
$s_a(v)=0$, which makes $\Psi_av=0$. After these conditions the form between $v$ and $e_r$ can be nonzero only when
$r$ is a row $(a,i)$ with $a\in\mathcal S_0$ and $i$ odd. So the radical contains at least
$\#(\mathcal S_0\cap\mathcal E)-\#(\mathcal S_1\setminus\mathcal E)$ independent vectors of $V$.

In both cases the vectors are supported on rows of integral weight. The $L_0$ rows of the zero class impose at most
$L_0$ further linear conditions. So at least $k_p$ independent vectors $u_1,\dots,u_{k_p}$ over $\F_p$ remain, with
\[
 \sum_s u_{j,s}\tilde G_{st}\equiv0\pmod p\quad\text{for every row $t$ of integral weight.}
\]

\begin{lemma}[weight raising]\label{lem:raise}
Let $E$ be a $\Z_{(p)}$-unimodular basis with half-integral weights $w$ for which the Hankel entries satisfy
$\vG(G_{st})\ge w_s+w_t$. Suppose that $k$ independent vectors over $\F_p$, supported on the rows of integral weight,
lie in the left radical of $\tilde G\bmod p$ against every row of integral weight. Then $\vG(\Delta_K)\ge2\sum_sw_s+k$.
\end{lemma}

\begin{proof}
Complete the lifts of the $k$ vectors to a unimodular change of basis that mixes only rows of integral weight. In the
new basis the $k$ retained rows have entries of valuation at least $w_s+w_t+1$ against the rows of integral weight. Against
a row of half-integral weight the bound $w_s+w_t$ is a half-integer, so the integral valuation is at least
$w_s+w_t+\frac12$. Hence each retained row has its weight raised by $\frac12$, and the determinant gains $k$.
\end{proof}

Theorem~\ref{thm:parity} follows from Lemma~\ref{lem:raise} with the vectors above.

\begin{remark}
In the inner range $K/p\le200$ the gain is about $p\,\gref(K/p)$ units, against the cost $pR(K/p)$ of Fauzan's
bound. We use the gain only for $K/p\le200$, where its integral against $x^{-3}$ is computed exactly.
\end{remark}

%% =====================================================================
\section{The prime sum}\label{sec:sum}

The prime number theorem converts the sums over $p$ into integrals. For a bounded piecewise continuous function $f$ on
$[c_0,c_1]\subset(0,\infty)$,
\begin{equation}\label{eq:pnt}
 \frac1K\sum_{c_0K<p\le c_1K}f(p/K)\log p\longrightarrow\int_{c_0}^{c_1}f(y)\,dy\qquad(K\to\infty).
\end{equation}
For crude bounds we use Chebyshev's inequality $\theta(y)=\sum_{p\le y}\log p\le y\log4$.

\begin{theorem}[denominator exponent]\label{thm:lambda}
Let $M(K)=\max(40,\lfloor K^{1/4}\rfloor)$ and let $m_K$ be the normalizing factor of Definition~\ref{def:LpC} with
$M=M(K)$. Then for all large $n$, $m_KF_K\in\Z[X]$ and
\[
 \log m_K\le\Lambda K^2,\qquad\Lambda=\frac{1288003}{10^6}.
\]
\end{theorem}

The theorem rests on the following integrals.

\begin{proposition}\label{prop:integrals}
\begin{enumerate}[label=(\alph*)]
\item $\displaystyle\int_{1/3}^{2\lambda}T_0(y)\,dy=\frac{127751}{96000}$.
\item $\displaystyle S_{\rm out}=\int_1^3\bigl[xT_0(1/x)-R(x)\bigr]_+\frac{dx}{x^3}=\frac{3203}{163200}$.
\item $\displaystyle\int_3^{200}R(x)\frac{dx}{x^3}=I_{200}$, an explicit rational number with
 $I_{200}=-0.01297481\ldots$
\item $\displaystyle\int_{200}^\infty R(x)\frac{dx}{x^3}\le-\frac{103109}{24000000}$.
\item $\displaystyle\int_3^{200}\gref(x)\frac{dx}{x^3}>\frac{5847}{10^6}$.
\end{enumerate}
\end{proposition}

\begin{proof}
Part (a) is Fauzan's \cite[(5.10)]{Fauzan2026}. In (a), (b), (c) and (e) the integrands are polynomials in $x$ and
$x^{-1}$ on finitely many intervals, on which the floor functions are constant and the minima and positive parts do not
change branch. The breakpoints are rational, so the integrals are rational numbers, computed exactly.

For (d) write $R=xF+Q$ by Lemma~\ref{lem:profile}, and put $f=\fr x$, $g=\fr{\alpha x}$, $G_0(v)=v^2/2-v^3/3-v/6$ and
$B_4(v)=v^2(1-v)^2-1/30$. Let
\begin{gather*}
 P=74g(1-g)-\lambda f(1-f),\qquad C=\tfrac{74}{\alpha}G_0(g)-\lambda G_0(f),\\
 D_3=-\frac{74B_4(g)}{12\alpha^2}+\frac{\lambda B_4(f)}{12}.
\end{gather*}
These functions are continuous and periodic in $f$ and $g$. Off $\Z\cup\alpha^{-1}\Z$ they satisfy $P'=F+\lambda$,
$C'=P-\bar P$ and $D_3'=C$, where $\bar P=2923/240$. Put $\Psi=\lambda/x+(P-\bar P)/x^2+2C/x^3+6D_3/x^4$. Then
$\Psi'=F/x^2-24D_3/x^5$. Since $D_3\le D_+=74/(360\alpha^2)+7\lambda/2880$ and $Q\le\frac{13}8$, the function
$\Omega=\Psi-6D_+/x^4-\frac{13}{16}x^{-2}$ satisfies $R(x)x^{-3}\le\Omega'(x)$ off a countable set. It is continuous,
and $\Omega(x)\to0$ as $x\to\infty$. So $\int_{200}^\infty R(x)x^{-3}dx\le-\Omega(200)$. At $x=200$ both $f$ and $g$
vanish, so $P=C=0$ and $D_3=74/(360\alpha^2)-\lambda/360$. Hence
\[
 \int_{200}^\infty\frac{R(x)}{x^3}\,dx\le-\frac\lambda{200}+\frac{\bar P+13/16}{200^2}+\frac{\lambda}{32\cdot200^4}
 \le-\frac{103109}{24000000}.
\]
\end{proof}

\begin{proof}[Proof of Theorem~\ref{thm:lambda}]
By Proposition~\ref{prop:normaliser}, $\log m_K=\sum_p(-L_p)\log p$. We bound each band.

\emph{Small primes}, $p\le K/M$. Here $-L_p\le6h\lfloor\log_p5K\rfloor+hv_p(24)$. The primes $p\le\sqrt{5K}$
contribute $O(K^{3/2}\log K)$. For larger $p$ the floor is $1$, and Chebyshev's inequality bounds the sum by
$7K\log4\cdot K/M$. Since $M(K)\to\infty$, the band contributes $o(K^2)$.

\emph{Inner primes}, $K/M<p\le K/3$. By \eqref{eq:innerunif} and Proposition~\ref{prop:kcount},
$-L_p\le pR(K/p)-[K\le200p]\,p\,\gref(K/p)+190M^3+30M$. The error sums to $O(M^3K)=O(K^{7/4})$. Put $x_1=10^8$.
For $K/x_1<p\le K/3$, \eqref{eq:pnt} in the variable $x=K/p$ turns the main terms into
$K^2\bigl(\int_3^{x_1}R(x)x^{-3}dx-\int_3^{200}\gref(x)x^{-3}dx\bigr)+o(K^2)$. By Lemma~\ref{lem:profile},
$|F|\le8\lambda$ and $|Q|\le\frac{13}8$, so $|R(x)|\le22x$ for $x\ge1$. So $p|R(K/p)|\le22(K+p)\le44K$ for $K/M<p\le K/x_1$, and Chebyshev's inequality bounds this
band by $(44\log4/x_1)K^2$.

\emph{Outer primes}, $K/3<p\le2h$. Here $-L_p\le\min\bigl(KT_0(p/K),\,pR(K/p)\bigr)+O(1)$ for $p\le K$, and
$-L_p=KT_0(p/K)+O(1)$ for $K<p\le2h$. By \eqref{eq:pnt} the sum is at most
$K^2\bigl(\int_{1/3}^{2\lambda}T_0-S_{\rm out}\bigr)+o(K^2)$.

Adding the bands and using Proposition~\ref{prop:integrals}, $\limsup K^{-2}\log m_K$ is at most
\begin{multline*}
 \frac{127751}{96000}-\frac{3203}{163200}+I_{200}-\frac{103109}{24000000}-\frac{5847}{10^6}\\
 +\frac{44\log4}{10^8}+\int_{10^8}^\infty\frac{|R(x)|}{x^3}\,dx .
\end{multline*}
The first five terms add up to a number in $(1.287995331,\,1.287995332)$. Since $|R(x)|\le22x$, the last two terms
are at most $44\log4/10^8+22/10^8<10^{-6}$, and the sum is below $\Lambda$.
\end{proof}

\begin{remark}
Fauzan's bound at the cutoff $M=100000$ is $A_{100000}=1.31757\ldots$ \cite[(5.20)]{Fauzan2026}. Of the gain of
$0.0296$, the outer extension gives $3203/163200\approx0.0196$ and the parity gain at least $0.0058$. The remaining
$0.0041$ comes from the exact integral to $200$, the sharper tail and the growing cutoff, which removes the terms of
order $1/M$. Computations of the exact content of $\Delta_K$ for
$K\le800$ suggest that the true exponent is near $1.2735$. They are not part of the proof.
\end{remark}

%% =====================================================================
\section{Decay}\label{sec:decay}

\subsection{The energy bound}
Fauzan bounds $\Delta_K(\xi)$ through \eqref{eq:andreief}. After the substitution $y_i=K\sqrt{t_i}$, the logarithm of
the integrand is, up to lower order terms, $K^2$ times a discrete logarithmic energy of the points $t_i$ in the field
$V$ of \eqref{eq:field}. We use the logarithmic potential and energy in Fauzan's sign convention,
\[
 U^\sigma(t)=\int\log|t-u|\,d\sigma(u),\qquad I(\sigma)=\iint\log|t-u|\,d\sigma(t)\,d\sigma(u).
\]
Let $\omega_{[0,b]}$ be the arcsine probability measure of $[0,b]$. We say that a positive measure $\sigma$ has an
\emph{arcsine majorant} if $\sigma\le\Theta\,\omega_{[0,b]}$ for a constant $\Theta$.

\begin{theorem}[energy bound]\label{thm:energy}
Let $0<b<4$ and let $\sigma$ be a positive measure of mass $\lambda$ on $[0,b]$ with an arcsine majorant. Let
$M_0\ge-8$ satisfy
\begin{equation}\label{eq:cellhyp}
 2U^\sigma(t)-V(t)\le M_0\qquad(0\le t\le4).
\end{equation}
Put $C^*=-2\lambda+12\alpha\lambda(1-\log\alpha)+3\lambda^2-2\lambda^2\log(2\lambda)$. Then for every $n\ge1$,
\[
 \log F_K(\xi)\le\bigl(C^*+\lambda M_0-I(\sigma)\bigr)K^2+24K\log K+O(K).
\]
\end{theorem}

The constant $C^*$ comes from Stirling's formula for $S_K$ and from the node sums, as in \cite[(6.3)]{Fauzan2026}. The
proof follows \cite[\S6]{Fauzan2026}, with three changes.
\begin{enumerate}[label=(\alph*),leftmargin=*]
\item \emph{The tail.} For $t\ge4$ the inequality $2U^\sigma-V+\frac12\sqrt t\le M_0$ holds for every such
 $\sigma$ and every $M_0\ge-8$. Indeed $U^\sigma(t)\le\lambda\log t$, $\int_0^1\log(t+u^2)du\ge\log t$ and
 $\int_0^\alpha\log(t+u^2)du\le\alpha\log t+\alpha^3/(3t)$. With $s=\sqrt t\ge2$ the resulting bound
 $1.3\log t+2\alpha^3/t-(2\pi-\frac12)s$ is at most $-8.96$.
\item \emph{The zero-mass lemma.} Let $\sigma_1,\dots,\sigma_r$ be finite measures of bounded support on $\C$ and
 $s_1,\dots,s_r$ real numbers with $\sum_ks_k\sigma_k(\C)=0$. If the logarithm is integrable, then
 $\sum_{k,l}s_ks_l\iint\log|z-w|\,d\sigma_k\,d\sigma_l\le0$. The proof truncates the kernel, writes it through
 Frullani's integral as a superposition of Gaussian kernels, and uses their positive definiteness.
\item \emph{The regularization.} The points $t_i$ are replaced by circles of radius $\epsilon$ and compared with
 $K\sigma$ through (b). The cross term of a circle and $\sigma$ is bounded by $U^\sigma(t_i)$ plus
 $\int\log^+(\epsilon/|t-u|)\,d\sigma(u)\le2\sqrt2\,\Theta\sqrt{\epsilon/b}$. This follows from the layer-cake formula
 and $\omega_{[0,b]}([x-r,x+r])\le\sqrt{2r/b}$. It also shows that $\sigma$ has no atoms and that $\log|t-u|$ is
 $\sigma$-integrable.
\end{enumerate}
With these, the configuration inequality gives, for distinct $t_i\ge0$,
\[
 2\sum_{i<j}\log|t_i-t_j|-K\sum_iV(t_i)+\sum_i\sqrt{t_i}\le(\lambda M_0-I(\sigma))K^2+O(K\log K).
\]
The weight bound $w(y)\le8192(1+y)^5e^{-2\pi y}$, the comparison of the products $D_N^6/D_K$ with $e^{-KV}$ and
Stirling's formula for $S_K$ then give Theorem~\ref{thm:energy}.

\subsection{The comparison measure}
Let $T_k$ be the Chebyshev polynomials and let
\[
 f(x)=\sum_{k=0}^{60}c_kT_k(x),\qquad \sigma=f\bigl(\tfrac{2t}b-1\bigr)\,\omega_{[0,b]}(dt),\qquad b=\Bigl(\frac{43}{50}\Bigr)^2.
\]

\begin{lemma}[closed forms]\label{lem:closed}
Suppose $f\ge0$ on $[-1,1]$. Then $\sigma$ has mass $c_0$ and $\sigma\le(\sum_k|c_k|)\,\omega_{[0,b]}$. With
$x=2t/b-1$,
\[
 U^\sigma(t)=c_0\log\frac b4-\sum_{k\ge1}\frac{c_kT_k(x)}k\qquad(0\le t\le b),
\]
\[
 U^\sigma(t)=c_0\Bigl(\log\frac b4+w\Bigr)-\sum_{k\ge1}\frac{c_ke^{-kw}}k\qquad\bigl(t=\tfrac b2(1+\cosh w),\ w\ge0\bigr),
\]
and
\[
 I(\sigma)=c_0^2\log\frac b4-\sum_{k\ge1}\frac{c_k^2}{2k}.
\]
\end{lemma}

\begin{proof}
For $y\in[-1,1]$ and $x=\cos\theta\in[-1,1]$, $\log|y-x|=-\log2-\sum_{k\ge1}\frac2kT_k(y)\cos k\theta$. For $y$
outside, $y=\cosh w$, the same expansion holds with $\cos k\theta$ replaced by $e^{-kw}\cos k\theta$ and $-\log2$ by
$w-\log2$. Integrating against $f(\cos\theta)\,d\theta/\pi$ and using orthogonality gives the potentials, after the
change of scale from $[-1,1]$ to $[0,b]$. Integrating the inner formula against $\sigma$ and using
$\int T_k\,d\sigma=c_k/2$ for $k\ge1$ gives the energy.
\end{proof}

\begin{lemma}[positivity]\label{lem:sos}
Let $h_0,\dots,h_{60}$ be real numbers, not all zero. Put $a_0=\sum_jh_j^2$, $a_k=2\sum_jh_jh_{j+k}$ for $k\ge1$,
and $c_k=(\lambda/a_0)a_k$. Then $f(\cos\theta)=(\lambda/a_0)\bigl|\sum_jh_je^{ij\theta}\bigr|^2\ge0$, and $c_0=\lambda$.
\end{lemma}

The proof is the expansion of the square. We take for $h_0,\dots,h_{60}$ the integers of the formal proof. They begin
$-2573,-2919,-6471,-5888,\dots$. These data were found by numerical optimization and play no other role in the proof.
Put $M_0=-3330031/500000=-6.660062$.

\subsection{The cells}
Condition \eqref{eq:cellhyp} is checked on two regions. On $[0,b]$ put $t=b\cos^2\varphi$ with
$0\le\varphi\le\pi/2$. By Lemma~\ref{lem:closed}, $2U^\sigma(t)-V(t)$ becomes
\[
 g_A(\varphi)=2\lambda\log\frac b4-2\sum_{k=1}^{60}\frac{c_k\cos2k\varphi}k-V(b\cos^2\varphi).
\]
On $[b,4]$ put $t=\frac b2(1+\cosh x)$ with $0\le x\le3$. This covers $[b,4]$, since $\frac b2(1+\cosh3)>4$. Then
\[
 g_B(x)=2\lambda\Bigl(\log\frac b4+x\Bigr)-2\sum_{k=1}^{60}\frac{c_ke^{-kx}}k-V\bigl(\tfrac b2(1+\cosh x)\bigr).
\]
In both regions the field is written in $y=\sqrt t$, with
$\int_0^A\log(y^2+u^2)\,du=A\log(y^2+A^2)-2A+2y\arctan(A/y)$. Its derivatives in $y$ have poles only at $y=\pm i$ and
$y=\pm i\alpha$.

\begin{proposition}[cell certificate]\label{prop:cells}
$g_A\le M_0$ on $[0,\pi/2]$ and $g_B\le M_0$ on $[0,3]$. Hence \eqref{eq:cellhyp} holds.
\end{proposition}

\begin{proof}
Region A is covered by $100$ cells $[\varphi_c-r,\varphi_c+r]$ with centres $\varphi_c=2\arctan m$, $m$ rational. Then
$\cos\varphi_c$ and $\sin\varphi_c$ are rational, and so are all $\cos2k\varphi_c$ and $\sin2k\varphi_c$, by the
rotation recurrence. On each cell the function is bounded by its Taylor polynomial of order four at the centre and a
Lagrange remainder of order five. The quadratic part is maximized exactly. The remainder of the trigonometric sum is
bounded by $2\sum_k|c_k|(2k)^5/k$. The remainder of the field is bounded by Fa\`a di Bruno's formula, with the
derivatives of $y=\sqrt b\cos\varphi$ and the distances to the poles of the field. The values of $V$ and its first
derivative at the centre need rational enclosures of $\log$ and $\arctan$ at rational points. The overlaps of
consecutive cells are checked with $\arctan m'-\arctan m\le(m'-m)/(1+mm')$. Region B is covered in the same way by
$10$ cells with centres $x_c=2\log r$, $r\ge1$ rational, so that $e^{-kx_c}$ and $\cosh(x_c/2)$ are rational. Every
cell inequality is a finite computation in rational arithmetic, carried out by the Lean kernel.
\end{proof}

\begin{theorem}[decay]\label{thm:decay}
$F_K(\xi)\le e^{-(D-o(1))K^2}$ with $D=138500086/10^8$.
\end{theorem}

\begin{proof}
By Lemma~\ref{lem:sos} and Lemma~\ref{lem:closed}, $\sigma$ is a positive measure of mass $\lambda$ on $[0,b]$ with an
arcsine majorant. By Proposition~\ref{prop:cells} it satisfies \eqref{eq:cellhyp}. Theorem~\ref{thm:energy} gives
$\limsup K^{-2}\log F_K(\xi)\le-(I(\sigma)-\lambda M_0-C^*)$. By Lemma~\ref{lem:closed},
$I(\sigma)-\lambda M_0-C^*=\lambda^2\log(b/4)-\sum_kc_k^2/(2k)-\lambda M_0-C^*=1.38500086079\ldots$. The inequality
with $D$ needs three rational enclosures of logarithms and exact arithmetic.
\end{proof}

\begin{remark}
Fauzan's measure, a sum of sixteen arcsine measures on nested intervals, gives $D=1.36699$
\cite[Lemma~6.1]{Fauzan2026}. The value $D$ of the method is the maximum of $I(\sigma)-\lambda M_0(\sigma)-C^*$ over
measures of mass $\lambda$, where $M_0(\sigma)$ is the supremum of $2U^\sigma-V$. It is attained at the equilibrium
measure of mass $\lambda$ in the field $V/2$ \cite{SaffTotik}. Numerical solution of this problem gives
$D\approx1.38551$. So the remaining room on the analytic side is about $5\cdot10^{-4}$. This value is a numerical
observation and is not part of the proof.
\end{remark}

%% =====================================================================
\section{Proof of Theorem~\ref{thm:main}}\label{sec:proof}

Let $m_K$ be the normalizing factor of Theorem~\ref{thm:lambda} and put $P_n=m_KF_K$, of degree $h=\lambda K$. Its
coefficients are integers, so condition (a) of Theorem~\ref{thm:criterion} holds with $d_n=1$ and the exponent $0$.
By Proposition~\ref{prop:rho}, $\rho=907427/10^5$ is a norm exponent. Let $\varepsilon>0$ and
$|t|\le e^{-(\rho+\varepsilon)K}$. Corollary~\ref{cor:nearby} gives $0<F_K(\xi+t)\le2F_K(\xi)$. With
Theorems~\ref{thm:lambda} and~\ref{thm:decay},
\[
 |P_n(\xi+t)|\le2e^{\Lambda K^2}e^{-(D-\varepsilon)K^2},
\]
so (b) holds with the rate $D-\Lambda$. Corollary~\ref{cor:nearby} and $m_K>0$ give (c). Theorem~\ref{thm:criterion}
applies with $\gamma=\lambda$, and
\[
 \frac{\lambda\rho}{D-\Lambda}=\frac{37}{40}\cdot\frac{9.07427}{1.38500086-1.288003}=86.5349\ldots<87 .
\]
Put $E=\lambda\rho/(D-\Lambda)=839369975/9699786$. Theorem~\ref{thm:criterion} shows that $\zeta(5)$ is irrational and
that the hypothesis of Lemma~\ref{lem:plain} holds with this $E$. Since $E<87$, Lemma~\ref{lem:plain} with $\mu_0=87$
gives an integer $q_0$ such that $|\zeta(5)-p/q|>q^{-87}$ for all integers $p$ and all integers $q\ge q_0$. This is
Theorem~\ref{thm:main}.

%% =====================================================================
\section{Formalization}\label{sec:lean}
\raggedbottom

Theorem~\ref{thm:main} is formalized in Lean~4 \cite{Lean4} with Mathlib \cite{Mathlib}, together with every theorem,
proposition and lemma its proof uses. The remarks are not part of the formalization. The development is in the
directory \lean{CatalanCore/Zeta5} of the repository
\begin{center}
\FormalizationLink
\end{center}
Theorem~\ref{thm:main} is the theorem
\begin{center}
{\small\texttt{Zeta5IrrationalityMeasure.zeta5\_irrationalityMeasure\_87}}
\end{center}
in \lean{SolutionZeta5.lean}. Its statement uses only Mathlib. The number $\zeta(5)$ is written as the series
$\sum_{n\ge0}(n+1)^{-5}$, and the conclusion is $q^{-87}<|\zeta(5)-p/q|$ for all integers $p$ and all sufficiently
large $q$. In the development $\zeta(5)$ is \lean{Apery.zeta5}, and \lean{riemannZeta_five} identifies it with
Mathlib's $\zeta(5)$. The root is proved in the module \lean{CatalanCore.Zeta5.Data.Public}. There
\lean{exponentE} is the exponent $E=839369975/9699786<87$, \lean{zeta5_exact} is the second statement of
Theorem~\ref{thm:criterion} at this exponent, \lean{irrational_of_family_nnv} is its first statement, and
\lean{TrigammaDifferences.Aux.eventually_of_point} is Lemma~\ref{lem:plain}. The modules contain no \lean{sorry} and no
\lean{native_decide}. The command \lean{\#print axioms} reports only \lean{propext}, \lean{Classical.choice} and
\lean{Quot.sound}. The proof is kernel-checked. It has not been run through the independent checker
\lean{lake comparator} \cite{Comparator}.

Fauzan's construction, the integral representation, positivity, the degree, Andr\'eief's formula and parts of the
energy argument are ported from Firsching's formalization \cite{Firsching26}, under the Apache License~2.0, at commit
\texttt{7fe7367}. The ported files are in \lean{CatalanCore/Zeta5/Port} and are listed in the repository's
\texttt{NOTICE}. The prime number theorem is taken from the PNT+ project \cite{PNTPlus}, through OpenAI's workspace
\cite{OpenAI2026}. The proofs were checked with Lean \texttt{v4.35.0-rc2} and the matching Mathlib.

The results of the paper correspond to the following declarations. Their files are under
\lean{CatalanCore/Zeta5}. We give the names with a short prefix only. The full namespaces are
\lean{CatalanCore.Zeta5} and its subnamespaces, \lean{CatalanCore.Uniform} and \lean{Apery}.
\begin{itemize}
\item Theorem~\ref{thm:criterion}: \lean{measure_of_family_nnv_period} (the period $40$ of the scale) and
 \lean{Data.Public.irrational_of_family_nnv}, applied in \lean{Data.Public.zeta5_exact} to the family $m_KF_K$.
 Lemma~\ref{lem:plain}: \lean{TrigammaDifferences.Aux.eventually_of_point}.
\item Lemma~\ref{lem:relnorm}: \lean{Apery.det_pencil_bounds}. Corollary~\ref{cor:nearby}: \lean{F_pencil_bounds},
 \lean{zeta5_nearbyNV} and \lean{zeta5_hasScaledDecay}.
\item Proposition~\ref{prop:basic}: \lean{Apery.moment_rep}, \lean{Apery.G_posDef}, \lean{Apery.natDegree_}$\Delta$, and
 \eqref{eq:andreief} is \lean{Apery.}$\Delta$\lean{_eq_integral}.
\item Proposition~\ref{prop:rho}: \lean{Rho.zeta5_rho_bound}, from \lean{Rho.residue_identity},
 \lean{Rho.rj_abs_le}, \lean{Rho.gram_lower}, \lean{Rho.node_bound}, \lean{Rho.Vf_le_max},
 \lean{Rho.Phi_tangent} and \lean{Rho.rho_numerics}.
\item Theorem~\ref{thm:outer}: \lean{zc18_ext} and \lean{zc19_gamma}. The cost estimates are \lean{zc17_gammaIn},
 \lean{zc19_cost} and \lean{zc20_outer}. Lemma~\ref{lem:profile} is \lean{neg_Gam_sub_N0} with
 \lean{Profile.Qc_ge} and \lean{Profile.Qc_le}.
\item Definition~\ref{def:LpC} and Proposition~\ref{prop:normaliser}: \lean{Content.LpC}, \lean{Content.mC},
 \lean{Content.log_mC}, \lean{Content.GV_F_LpC} and \lean{Content.zc32_normaliser}.
\item Theorem~\ref{thm:parity}: \lean{Apery.parity_gain}, a proof of the statement \lean{Apery.ParityGain}. Its parts
 are \lean{tauX_lead} (Lemma~\ref{lem:taulead}), \lean{fullSym_odd} and \lean{locSymb_pc} (the reflection),
 \lean{lead_entry} (the identification \eqref{eq:leadform}), \lean{radical_T_odd} and \lean{radical_T_even} (the
 radical vectors), \lean{cclass_unimodular} and \lean{centred_entry} (the centred basis), and
 \lean{parity_gain_delta} (Lemma~\ref{lem:raise}). Proposition~\ref{prop:kcount} is \lean{Apery.kpar_ge}, with
 $0\le\gref\le\frac12$ in \lean{Profile.gref_nonneg} and \lean{Profile.gref_le_half}.
\item Proposition~\ref{prop:integrals}: \lean{Sum.zc34}, \lean{Sum.zc35}, \lean{Sum.zc36}, \lean{Sum.zc37} and
 \lean{Sum.zc38}. Their sum is \lean{Sum.lambdaQ_value}.
\item Theorem~\ref{thm:lambda}: \lean{Content.zc32_signed}, through \lean{Content.log_mC_le}, \lean{Content.pt_bound}
 and the prime sums \lean{Sum.weighted_prime_sum_tendsto} and \lean{Sum.theta_le_linear}. The bound $|R(x)|\le22x$ is
 \lean{Profile.Rg_bound}, and the band bound $(44\log4/x_1)K^2$ is the second term in the statement of \lean{Content.log_mC_le}. The tail
 in Proposition~\ref{prop:integrals}(d) is \lean{Sum.TL200}.
\item Theorem~\ref{thm:energy}: \lean{Energy.real_bound}, with the zero-mass lemma \lean{Energy.energy_log_nonpos},
 the regularization \lean{Energy.integral_regTerm_le}, the configuration inequality \lean{Energy.config_ineq} and the
 tail \lean{Decay.cellHyp_tail}.
\item Lemmas~\ref{lem:closed} and~\ref{lem:sos}: \lean{Potential.logPot_nu_inside}, \lean{Potential.logPot_nu_cosh},
 \lean{Potential.energy_nu}, \lean{Potential.chebF_sos} and \lean{Potential.sup_le_sumAbs}.
\item Proposition~\ref{prop:cells}: \lean{Cells.cellA_sound}, \lean{Cells.cellB_sound} and \lean{Cells.hcell_of_cells},
 applied to the $100$ and $10$ cells of \lean{Data/CheckA1.lean}--\lean{CheckA4.lean} and \lean{Data/CheckB.lean}.
\item Theorem~\ref{thm:decay}: \lean{Decay.zeta5_decay}, with the constant inequality \lean{Cells.zd22_sound}.
\end{itemize}
The certificate data are the $61$ integers $h_j$, the cells, and the rational breakpoints of the integrals of
Proposition~\ref{prop:integrals}. They were computed with floating-point optimization and the interval arithmetic of
Arb \cite{Arb}, and Lean re-checks them with rational arithmetic. They play no role in the proof beyond satisfying the
checks.

\section*{Disclosure of AI use}
The mathematics, the Lean formalization, the certificate generators and the text of this paper were produced by AI
systems at the request of the author. The author posed the problems, directed the work and is responsible for the
paper. The systems were Claude Opus~5.5, Claude Sonnet~5.5 and Claude Haiku~5.5 (Anthropic), running in Claude Code,
and GPT-6.1 Sol (OpenAI), running in Codex. The construction and the irrationality proof of $\zeta(5)$ are due to
Fauzan \cite{Fauzan2026}.

\subsection*{Acknowledgements}
We thank the authors of Mathlib and of the PNT+ project, Moritz Firsching for his formalization of Fauzan's proof, and
OpenAI for making its formalization available.

\bibliographystyle{amsplain}
\bibliography{refs}

\end{document}
